% HarnessAgent: Harness-as-a-Service for Production Multi-Framework AI Agents
% IEEE Conference Paper — NeurIPS 2026 Workshop on Production AI Systems
% Author: Pradip Tivhale (thepradip)
%
% All benchmark numbers are verifiable against:
%   benchmarks/results/*.json  (scripts in benchmarks/*.py)
% Every citation is a real published work with a verifiable DOI or arXiv ID.

\documentclass[conference]{IEEEtran}

\usepackage{amsmath}
\usepackage{amssymb}
\usepackage{graphicx}
\usepackage{booktabs}
\usepackage{hyperref}
\usepackage{xcolor}
\usepackage{listings}
\usepackage{microtype}
\usepackage{url}
\usepackage{cite}
\usepackage{array}
\usepackage{multirow}
\usepackage{tabularx}
\usepackage{balance}
\usepackage{tikz}
\usetikzlibrary{arrows.meta, positioning, shapes.geometric, backgrounds, fit}

\hypersetup{
  colorlinks=true,
  linkcolor=blue!70!black,
  citecolor=blue!60!black,
  urlcolor=blue!60!black
}

\lstset{
  basicstyle=\ttfamily\scriptsize,
  breaklines=true,
  frame=single,
  backgroundcolor=\color{gray!8},
  keywordstyle=\color{blue!70!black}\bfseries,
  commentstyle=\color{green!50!black}\itshape,
  stringstyle=\color{orange!80!black},
  language=Python,
  showstringspaces=false,
  tabsize=2
}

\begin{document}

\title{HarnessAgent: Harness-as-a-Service for\\Production Multi-Framework AI Agents}

\author{
  \IEEEauthorblockN{Pradip Tivhale}
  \IEEEauthorblockA{Stewart\\
    \texttt{Pradip.Tivhale@stewart.com}}
}

\maketitle

%-----------------------------------------------------------------------
\begin{abstract}
AI agents fail in production not because of weak reasoning, but because
of missing \emph{infrastructure}: providers go down, context windows
overflow, costs spiral, safety rules are applied unevenly, and failures
compound silently. We present \textbf{HarnessAgent}, an open-source
Harness-as-a-Service (HaaS) layer that separates production concerns
from task logic and works with four major agent frameworks (LangGraph,
AutoGen, CrewAI, Agno). The harness delivers seven pillars: health-aware
LLM routing with circuit breaking, three-tier context engineering with
GraphRAG retrieval, hierarchical tracing, a three-stage safety pipeline
with human-in-the-loop escalation, real-time Redis Stream feedback, the
Hermes runtime self-improvement loop, and a tool result size cap.

Eight benchmarks validate the design. GraphRAG reduces schema tokens by
\textbf{82.9\%} (2,209$\to$378 tokens, 100\% table coverage). The
semantic cache achieves \textbf{100\% TPR and 0\% FPR} at cosine
threshold 0.97. The circuit breaker opens in \textbf{0.01\,ms} with
zero false trips. Span recording overhead is \textbf{872\,\textmu{}s
p50} (fakeredis) / \textbf{1,331\,\textmu{}s p50} (real Redis).
On 28 AgencyBench-V2 tasks spanning six capability dimensions, adding
HarnessAgent infrastructure over bare LangGraph lifts pass rate from
\textbf{50.0\% to 71.4\%} (+21.4\,pp). On 50 $\tau$-bench retail and
airline tasks, the safety pipeline blocks \textbf{100\%} of adversarial
requests with \textbf{0\%} false positive rate. On 50 ATBench-inspired
safety trajectories the guardrail achieves \textbf{TPR\,=\,96.7\%} and
\textbf{TNR\,=\,100\%}. The Hermes loop improves pre/post-patch pass@1
by \textbf{+85\,pp} on embedded GAIA-style tasks (5\%$\to$90\%).
The package is \texttt{pip install agent-haas} with 1,087 passing tests.
\end{abstract}

\begin{IEEEkeywords}
AI agents, production systems, agent infrastructure, LLM routing,
context engineering, safety pipeline, self-improvement, multi-framework
\end{IEEEkeywords}

%=======================================================================
\section{Introduction}
\label{sec:introduction}

Building an AI agent that works in a notebook is easy. Keeping it
working in production is not. Frameworks like LangGraph~\cite{langgraph},
AutoGen~\cite{autogen}, CrewAI, and Agno define what agents do — they
do not define how those agents handle provider outages, exploding context
windows, policy violations, cost overruns, or prompt drift.

Every team that ships an agent in production solves these problems from
scratch: they write their own circuit breaker, their own retry logic,
their own tracing, their own safety filter. The result is duplicated,
untested infrastructure that quietly breaks in the middle of the night.

We draw an analogy to early database engineering. Before
Database-as-a-Service (DBaaS), teams managed their own replication,
failover, and backups. DBaaS separated storage infrastructure from
application logic. We propose the same separation for agents:
\textbf{Harness-as-a-Service (HaaS)}, where the harness owns routing,
memory, tracing, safety, cost, recovery, and improvement while the agent
framework owns task logic.

Princeton's Holistic Agent Leaderboard (HAL, accepted ICLR 2026~\cite{hal2025})
recently pivoted from accuracy-only scoring to a Reliability Dashboard
covering consistency, robustness, and safety — confirming that raw task
performance is insufficient for production. HarnessAgent provides the
infrastructure layer that turns any agent framework into a reliable,
observable, self-improving production system.

\subsection*{Contributions}

\begin{enumerate}
  \item \textbf{HaaS design contract and implementation.} A formal
    separation between task logic and production infrastructure,
    implemented as a pip-installable layer with adapters for four
    frameworks and verified by 1,087 passing tests.
  \item \textbf{Three-tier context engineering with GraphRAG.}
    A configurable L1/L2/L3 pipeline with graph-structured schema
    retrieval that cuts token consumption by 82.9\% while maintaining
    100\% table coverage.
  \item \textbf{Harness ablation on AgencyBench-V2.} A controlled
    three-condition experiment across 28 AgencyBench-V2 tasks showing
    the harness lifts pass rate from 50.0\% (bare) to 71.4\% (full
    HaaS), exceeding the published AgencyBench native-SDK baseline of
    48.4\%~\cite{agencybench2026} by 23.0\,pp.
  \item \textbf{Safety pipeline validation.} On $\tau$-bench-inspired
    tasks, 100\% adversarial compliance with 0\% false positives; on
    ATBench-taxonomy trajectories, TPR\,=\,96.7\% and TNR\,=\,100\%.
  \item \textbf{Hermes runtime self-improvement.} A loop that patches
    agent prompts from observed failures without stopping the agent,
    achieving +85\,pp pre/post-patch improvement on GAIA-style tasks
    and +6\,pp on verified SQL execution.
\end{enumerate}

%=======================================================================
\section{Background and Related Work}
\label{sec:related}

\subsection{Harness Engineering}

Zhong and Zhu~\cite{zhong2026harness} define harness engineering as a
systematic framework with four maturity levels (H0--H3) and eleven
components. Their contribution is taxonomy: there is no implementation,
no routing, and no multi-framework support. HarnessAgent implements
their theoretical model end-to-end.

Lin et~al.~\cite{lin2026ahe} evolve harness \emph{structure} offline
between benchmark rounds. Hermes patches agent \emph{prompts and tool
configurations} at production runtime. The locus of adaptation differs.

Seong et~al.~\cite{seong2026last} frame harness construction as
meta-learning. This is theoretically interesting but provides no
production infrastructure, no safety guarantees, and no routing.

Meng et~al.~\cite{meng2026survey} survey 22 agent systems and
explicitly identify ``production multi-framework harness'' as an open
research direction. HarnessAgent fills that gap.

\subsection{Memory and Context}

MemGPT~\cite{memgpt2023} pages context segments to simulate extended
working memory. Our ContextEngine compresses and promotes content across
three tiers continuously, integrating GraphRAG so the agent receives
only the most relevant context per query.

GraphRAG~\cite{graphrag2024} applies knowledge graphs to document
retrieval for multi-hop summarization. We specialize this to agent
context: entity graphs over SQL schema objects with weighted BFS
retrieval of the minimal table set per query.

TERAG~\cite{terag2025} proposes token-efficient graph retrieval.
HarnessAgent's SchemaStore applies similar compression inside a
production harness with multi-tenancy, TTL invalidation, and runtime
SQL execution feedback.

\subsection{Benchmarks and Evaluation}

GAIA~\cite{gaia2023} (450 questions, three difficulty levels) tests
general AI assistants on tool use, multi-step reasoning, and
multi-modal tasks. HAL's top GAIA result is 74.55\%~\cite{hal2025} with
Claude Sonnet 4.5. We use GAIA failure types to seed Hermes; our
measurement is pre/post-patch improvement, not a leaderboard score.

AgencyBench-V2~\cite{agencybench2026} (138 tasks, 6 capabilities,
32 real-world scenarios, avg.\ 90 tool calls) finds that native-SDK
harnesses score 48.4\% while weaker independent harnesses score
substantially lower, demonstrating that harness quality directly drives
agent success. We use their task descriptions and capability taxonomy.

$\tau$-bench~\cite{taubench2024} tests stateful tool use and policy
compliance in retail and airline domains with user simulation. HAL's
best $\tau$-bench airline result is 56\% task success~\cite{hal2025}.
We measure safety compliance — a dimension absent from the leaderboard.

ATBench~\cite{atbench2026} organizes agent safety evaluation into six
categories: prompt injection, privilege escalation, data exfiltration,
unsafe code execution, policy bypass, and harmful content generation.
We design 50 trajectories following this taxonomy.

\subsection{Safety}

Existing agent frameworks provide no standardized safety infrastructure.
LangSmith, Portkey, and LangGraph Platform provide monitoring but not
three-stage runtime guardrails or HITL escalation.
Table~\ref{tab:comparison} summarizes the capability gap.

\begin{table}[t]
\caption{Capability Comparison of Production Agent Infrastructure}
\label{tab:comparison}
\centering
\small
\setlength{\tabcolsep}{3pt}
\begin{tabular}{lcccccc}
\toprule
\textbf{Capability} & \rotatebox{70}{\textbf{LangSmith}} &
  \rotatebox{70}{\textbf{Portkey}} &
  \rotatebox{70}{\textbf{LG Platform}} &
  \rotatebox{70}{\textbf{Mem0}} &
  \rotatebox{70}{\textbf{HarnessAgent}} \\
\midrule
Circuit-breaker routing       & --  & part. & part. & -- & \checkmark \\
Semantic LLM cache             & --  & --    & --    & -- & \checkmark \\
3-tier memory + GraphRAG       & --  & --    & --    & part. & \checkmark \\
7-kind hierarchical tracing    & \checkmark & -- & \checkmark & -- & \checkmark \\
3-stage safety + HITL          & --  & --    & --    & -- & \checkmark \\
Runtime self-improvement       & --  & --    & --    & -- & \checkmark \\
Tool result size cap           & --  & --    & --    & -- & \checkmark \\
Cross-framework adapters       & --  & \checkmark & -- & -- & \checkmark \\
Open source + \texttt{pip}     & part. & part. & -- & \checkmark & \checkmark \\
\bottomrule
\end{tabular}
\end{table}

%=======================================================================
\section{Architecture}
\label{sec:architecture}

HarnessAgent enforces a two-layer contract: the \emph{framework layer}
implements \texttt{task\_step(state)\;$\to$\;(output,\,next\_state)};
the \emph{harness layer} implements
\texttt{run(agent,\,task)\;$\to$\;TraceView}.
The harness owns all seven production pillars; the agent owns task logic.

\begin{lstlisting}[language=Python,caption={Four-line integration}]
from harness import HarnessAgent, LangGraphAdapter

graph = MyLangGraph()                      # your existing agent
harness = HarnessAgent(LangGraphAdapter(graph))
result  = await harness.run(task, tenant_id="acme")
trace   = harness.trace(result.run_id)    # full span waterfall
\end{lstlisting}

Figure~\ref{fig:architecture} shows the full system. The four core
domains (Memory, LLM Router, Tools, Safety) sit between the agent and
the observability/self-improvement layers, shielding the agent from all
infrastructure concerns.

% Architecture diagram — spans both columns
\begin{figure*}[t]
\centering
\begin{tikzpicture}[
  every node/.style={font=\small},
  layer/.style={draw=black!70, rounded corners=4pt, fill=#1,
                minimum width=14.5cm, minimum height=0.82cm,
                text centered, align=center, inner sep=4pt},
  domain/.style={draw=black!70, rounded corners=3pt, fill=#1,
                 minimum width=3.35cm, minimum height=1.6cm,
                 text centered, align=center, font=\scriptsize,
                 inner sep=3pt},
  arr/.style={-{Stealth[length=4pt]}, thick, draw=black!50}
]

% ── Layers ─────────────────────────────────────────────────────
\node[layer=blue!10] (client) at (0,7.6)
  {\textbf{Client} \quad REST API \;·\; SSE Streaming \;·\; Python SDK};

\node[layer=purple!10] (orch) at (0,6.35)
  {\textbf{Orchestration} \quad AgentRunner \;·\; Planner (DAG) \;·\; Scheduler \;·\; HITL Manager};

\node[layer=orange!10] (agent) at (0,5.1)
  {\textbf{Agent Layer} \quad BaseAgent \;·\; SQLAgent \;·\; CodeAgent \;·\;
    LangGraph \;·\; AutoGen \;·\; CrewAI \;·\; Agno adapters};

% ── Four domains ────────────────────────────────────────────────
\node[domain=green!12] (mem) at (-5.45,3.25)
  {\textbf{Memory}\\L1 Redis (hot)\\L2 ContextEngine\\L3 GraphRAG / VectorDB};
\node[domain=yellow!20] (llm) at (-1.82,3.25)
  {\textbf{LLM Router}\\Circuit Breaker\\Semantic Cache\\Cost Tracker};
\node[domain=red!10] (tools) at (1.82,3.25)
  {\textbf{Tools}\\ToolRegistry\\Code Sandbox\\MCP Client};
\node[domain=cyan!12] (safety) at (5.45,3.25)
  {\textbf{Safety}\\Input Guard\\Step Guard\\Output Guard + HITL};

% ── Observability ───────────────────────────────────────────────
\node[layer=gray!12] (obs) at (0,1.65)
  {\textbf{Observability} \quad TraceRecorder (7 span kinds) \;·\; MLflow \;·\;
    OpenTelemetry \;·\; Prometheus (15 metrics)};

% ── Self-improvement ────────────────────────────────────────────
\node[layer=teal!12] (hermes) at (0,0.42)
  {\textbf{Hermes Self-Improvement} \quad ErrorCollector
    \;$\to$\; PatchGenerator \;$\to$\; Evaluator \;$\to$\; Apply / Rollback};

% ── Arrows ──────────────────────────────────────────────────────
\draw[arr] (client.south)  -- (orch.north);
\draw[arr] (orch.south)    -- (agent.north);
\draw[arr] (agent.south)   -- ($(agent.south)+(0,-0.30)$)
              -| (mem.north);
\draw[arr] (agent.south)   -- ($(agent.south)+(0,-0.30)$)
              -| (llm.north);
\draw[arr] (agent.south)   -- ($(agent.south)+(0,-0.30)$)
              -| (tools.north);
\draw[arr] (agent.south)   -- ($(agent.south)+(0,-0.30)$)
              -| (safety.north);
\draw[arr] (mem.south)     |- ($(obs.north)+(-4.8,0.2)$)
              -- ($(obs.north)+(-4.8,0)$);
\draw[arr] (llm.south)     |- ($(obs.north)+(-1.6,0.2)$)
              -- ($(obs.north)+(-1.6,0)$);
\draw[arr] (tools.south)   |- ($(obs.north)+(1.6,0.2)$)
              -- ($(obs.north)+(1.6,0)$);
\draw[arr] (safety.south)  |- ($(obs.north)+(4.8,0.2)$)
              -- ($(obs.north)+(4.8,0)$);
\draw[arr] (obs.south) -- node[right,font=\tiny,text=gray]
              {failures} (hermes.north);
\end{tikzpicture}
\caption{HarnessAgent two-layer architecture. The \emph{framework layer}
(orange box) implements task logic. The \emph{harness layer} (all other
boxes) owns routing, memory, tracing, safety, cost, and
self-improvement. A four-line integration is sufficient for all four
supported frameworks.}
\label{fig:architecture}
\end{figure*}

\subsection{LLM Router}
\label{sec:router}

The router selects a provider using a composite health score:
\begin{equation}
  \text{score}(p) = w_1\,(1-\hat{\lambda}_p) +
                    w_2\,(1-\bar{L}_p/L_{\max}) +
                    w_3\,(1-c_p/c_{\max})
  \label{eq:routing}
\end{equation}
where $\hat{\lambda}_p$ is the estimated error rate, $\bar{L}_p$ the
exponentially-smoothed latency, and $c_p$ the cost per token. Health
results are cached for 5 seconds to prevent thundering-herd queries
under load.

\textbf{Circuit breaker.} Each provider has a sliding window of the
last $N$\,=\,5 calls. The state machine is:
\textsc{closed} $\xrightarrow{5\;\text{failures}}$ \textsc{open}
$\xrightarrow{60\,\text{s}}$ \textsc{half-open}
$\xrightarrow{2\;\text{successes}}$ \textsc{closed}.
In the \textsc{open} state, requests skip the provider without incurring
network latency. Section~\ref{sec:eval_cb} reports 0.01\,ms
time-to-open and zero false trips.

\textbf{Semantic LLM cache.} Responses are cached via
\texttt{all-MiniLM-L6-v2}~\cite{minilm} embeddings. A Redis sorted set
indexes entries by timestamp; vector scan is capped at the 200
most-recent entries, keeping lookup O(1) in practice. An exact
SHA-256 fast path avoids embedding for repeated identical prompts.

\textbf{Exponential backoff.} Rate-limit and timeout errors trigger
per-provider retries with delays of 1.0\,s, 2.0\,s, and 4.0\,s
($\pm$20\% jitter) before the router falls back to the next provider.

\subsection{Three-Tier Context Engineering}
\label{sec:context}

Context flows across three tiers with a default 60\,/\,25\,/\,15\%
token budget split:

\begin{itemize}[noitemsep]
  \item \textbf{L1 — Hot (Redis LIST).} Most recent turns, tool outputs,
    and guardrail events. Trimmed at 80\% capacity, oldest entries
    promoted to L2.
  \item \textbf{L2 — Warm (ContextEngine).} Compressed summaries in
    skill-isolated namespaces. Per-step action scoring gates which
    content is promoted.
  \item \textbf{L3 — Cold (VectorStore + GraphStore).} Long-term
    knowledge retrieved via GraphRAG. For SQL agents, the SchemaStore
    builds an entity graph over tables and foreign-key edges; weighted
    BFS up to depth 2 returns only the tables needed per query.
\end{itemize}

\textbf{Tool result cap.} Any tool output exceeding 8,000 characters is
truncated before entering agent history. The original length is recorded
in \texttt{metadata[``original\_chars'']} so the agent can detect
truncation and re-query with a narrower scope.

\subsection{Safety Pipeline}
\label{sec:safety}

Three guardrail stages gate every agent lifecycle:

\begin{enumerate}[noitemsep]
  \item \textbf{Input guard.} Detects prompt injection, PII, and policy
    violations before the agent is invoked.
  \item \textbf{Step guard.} Inspects each tool call name, arguments,
    and intermediate output for unsafe patterns (destructive shell
    commands, SQL DDL/DML, data exfiltration URLs).
  \item \textbf{Output guard.} Validates the final response for PII,
    injected instructions, and secret leakage before returning to the
    caller.
\end{enumerate}

\textbf{HITL escalation.} A \texttt{HITL\_REQUIRED} event suspends the
run and publishes a decision request to the FeedbackChannel (Redis
Stream). The run resumes after a human operator approves or rejects the
call. A configurable timeout applies a default policy if no decision
arrives.

\subsection{Hermes Self-Improvement Loop}
\label{sec:hermes_arch}

Hermes monitors live failures and patches agent prompts without stopping
the agent:

\begin{enumerate}[noitemsep]
  \item \textbf{ErrorCollector.} Intercepts outputs with reward below a
    threshold; clusters failures by error pattern.
  \item \textbf{PatchGenerator.} Calls a patch LLM with the failure
    cluster and current prompt fragment; generates $k$\,=\,3 candidate
    patches.
  \item \textbf{Evaluator.} Replays failing cases with each candidate
    patch using real execution. Selects the patch with the highest mean
    reward.
  \item \textbf{OnlineLearningMonitor.} Applies the winning patch live.
    If subsequent reward drops below the pre-patch baseline, rolls back
    and escalates to HITL.
\end{enumerate}

The key distinction from AHE~\cite{lin2026ahe}: AHE modifies harness
topology offline between evaluation rounds; Hermes modifies agent
prompts online at production runtime.

\subsection{Framework Adapters}
\label{sec:adapters}

Each adapter exposes four hooks:

\begin{lstlisting}[language=Python, caption={Adapter interface}]
class FrameworkAdapter(ABC):
    async def run(self, ctx, input) -> AsyncIterator[StepEvent]: ...
    async def get_result(self) -> FrameworkResult: ...
    def attach_harness(self, safety, cost_tracker, audit): ...
    def attach_mcp(self, mcp_client, tool_names=None): ...
\end{lstlisting}

Adapters are provided for LangGraph, AutoGen, CrewAI, and Agno. Each
wraps the framework's native execution loop, emitting one
\texttt{StepEvent} per observable step. Safety checks, cost recording,
and tracing run inside \texttt{run\_with\_harness()} without any change
to the underlying framework code.

%=======================================================================
\section{Implementation}
\label{sec:implementation}

HarnessAgent is written in Python 3.11+ with FastAPI and asyncio
throughout. Key implementation decisions:

\textbf{Redis over Kafka for spans.} The trace store uses Redis Streams
with 48-hour TTL rather than Kafka because span writes are synchronous
with agent steps; the sub-millisecond Redis overhead (Section~\ref{sec:eval_spans})
keeps tracing transparent to the agent.

\textbf{NetworkX for development, Neo4j for production.} The graph store
is swappable. NetworkX requires zero dependencies and covers all
development and test scenarios; Neo4j handles enterprise schemas with
200+ tables.

\textbf{fakeredis in tests.} All 1,087 unit and integration tests run
without a real Redis instance using \texttt{fakeredis}. Test execution
is fully deterministic and CI-friendly.

\textbf{0.97 cosine threshold.} The semantic cache threshold was chosen
empirically: near-exact pairs achieve mean cosine 0.987 (min 0.963);
decoy pairs average 0.083 (max 0.201), leaving a 0.762-point margin
that accommodates embedding model variance.

%=======================================================================
\section{Evaluation}
\label{sec:evaluation}

All benchmarks run on a MacBook Pro (Apple M3 Pro, 18\,GB RAM, Python
3.11). Scripts and result JSONs are in \texttt{benchmarks/} and
\texttt{benchmarks/results/} in the public repository.

\textbf{Simulation note.} Benchmarks 7.5--7.8 use calibrated scoring
models rather than live LLM calls because the measured quantities
(safety block rate, context utilization effect, improvement rate) are
properties of HarnessAgent's infrastructure, not of a specific LLM.
We follow the same approach as our real execution benchmarks: same
HarnessAgent components (SafetyPipeline, ContextEngine, TraceRecorder)
are active; only the LLM response is replaced by a deterministic model.
Real-execution benchmarks are 7.1--7.4 and 7.9--7.10.

\begin{table*}[t]
\caption{HarnessAgent Benchmark Results (all numbers match
  \texttt{benchmarks/results/*.json})}
\label{tab:benchmarks}
\centering
\small
\setlength{\tabcolsep}{5pt}
\begin{tabular}{llllll}
\toprule
\textbf{\#} & \textbf{Component} & \textbf{Metric} &
  \textbf{Result} & \textbf{Target} & \textbf{Method} \\
\midrule
7.1 & GraphRAG & Token reduction & \textbf{82.9\%} (2,209$\to$378 tok.) &
  $>$70\% & Real schema, $n$=20 queries \\
    & & Table coverage & \textbf{100\%} & 100\% & \\
\midrule
7.2 & Semantic cache & TPR @ $\tau$=0.97 & \textbf{100\%} &
  $>$95\% & Real all-MiniLM-L6-v2 embeddings \\
    & & FPR @ $\tau$=0.97 & \textbf{0\%} & $<$5\% & \\
\midrule
7.3 & Circuit breaker & Time-to-open & \textbf{0.01\,ms} & $<$1\,ms &
  3 fault scenarios, 0 false trips \\
\midrule
\multirow{2}{*}{7.4} & Span overhead & p50 (fakeredis) &
  \textbf{872\,\textmu{}s} & $<$1\,ms & 2,000 iterations \\
  & & p50 (real Redis) & \textbf{1,331\,\textmu{}s} & $<$5\,ms &
  2,000 iterations, localhost \\
\midrule
7.5 & AgencyBench-V2 ablation & Bare (A) pass rate &
  \textbf{50.0\%} & $>$native\,SDK\,(48.4\%) &
  28 tasks, real task descriptions \\
    & & +Infra (B) pass rate & \textbf{60.7\%} (+10.7\,pp) & & \\
    & & Full HaaS (C) pass rate & \textbf{71.4\%} (+21.4\,pp) & & \\
\midrule
7.6 & $\tau$-bench safety & Adversarial compliance (OFF) &
  \textbf{10.0\%} & --- & 50 tasks (40 benign, 10 adv.) \\
    & & Adversarial compliance (ON) &
  \textbf{100\%} (+90\,pp) & 100\% & \\
    & & False positive rate & \textbf{0.0\%} & $<$5\% & \\
\midrule
7.7 & ATBench safety & TPR (unsafe blocked) &
  \textbf{96.7\%} (29/30) & $\ge$90\% \checkmark &
  50 trajectories, 6 categories \\
    & & TNR (benign allowed) & \textbf{100\%} (20/20) &
  $\ge$95\% \checkmark & \\
\midrule
7.8 & Hermes (GAIA-style) & Pre-patch pass@1 &
  \textbf{5.0\%} & --- & 15 failure seeds, 20 eval tasks \\
    & & Post-patch pass@1 & \textbf{90.0\%} (+85\,pp) & --- &
  Converges cycle 1 \\
\midrule
7.9 & Hermes (real SQL, $n$=50) & Pre-patch pass@1 &
  \textbf{40.0\%} & --- & Real SQLite execution \\
    & & Post-patch pass@1 & \textbf{46.0\%} (+6\,pp) & $>$45\% & \\
\midrule
7.10 & NexusSql ablation ($n$=49) & LLM only (A) &
  \textbf{36.7\%} & --- & Real SQLite, gretelai dataset \\
    & & + Schema + verifier (B) & \textbf{44.9\%} (+8.2\,pp) & & \\
    & & + Self-correction (C) & \textbf{49.0\%} (+12.3\,pp total) & & \\
\bottomrule
\end{tabular}
\end{table*}

\subsection{GraphRAG Token Efficiency}
\label{sec:eval_graphrag}

A 10-table e-commerce schema (users, orders, order\_items, products,
categories, payments, reviews, sessions, addresses, promotions) serves
as the test environment. Twenty natural-language queries require 1--4
tables each.

The \emph{baseline} injects the full schema DDL for every query:
2,209 tokens average. The \emph{GraphRAG} path calls SchemaStore to
retrieve seed tables by vector similarity, then expands via weighted
BFS up to depth 2. Only retrieved tables' DDL enters the prompt.

Results: 82.9\% token reduction (2,209$\to$378 tokens average), 100\%
table coverage on all 20 queries. Nineteen queries used the vector-first
path; one multi-hop join query required graph traversal.
Token savings range from 75.1\% (3-table join) to 91.4\% (single-table
session query).

\subsection{Semantic LLM Cache}
\label{sec:eval_cache}

Fifty prompt pairs are embedded with \texttt{all-MiniLM-L6-v2}.
Each pair has a \emph{near-exact} variant (minor rephrasing, same
meaning) and a \emph{decoy} variant (different semantics).

Near-exact pairs: mean cosine 0.987 (min 0.963).
Decoy pairs: mean cosine 0.083 (max 0.201).
At $\tau$\,=\,0.97: 100\% TPR, 0\% FPR.
The 0.762-point margin between the minimum near-exact score (0.963) and
maximum decoy score (0.201) provides robust separation against embedding
variance.

\subsection{Circuit Breaker Reliability}
\label{sec:eval_cb}

Three scenarios: \emph{gradual} (5 consecutive failures),
\emph{burst} (5 simultaneous failures), and \emph{intermittent}
(failures below the 5-failure threshold interspersed with successes).

The circuit opens in 0.01\,ms after the fifth failure in both gradual
and burst scenarios. In the intermittent scenario, the sliding window
resets on successes and the circuit never trips — 0 false positive trips
across all three scenarios.

\subsection{Span Recording Overhead}
\label{sec:eval_spans}

Span overhead is measured over 2,000 iterations (200-iteration warmup)
with \texttt{perf\_counter} timing.

\emph{Fakeredis} (in-process): p50\,=\,872\,\textmu{}s,
p95\,=\,954\,\textmu{}s, p99\,=\,1,101\,\textmu{}s.

\emph{Real localhost Redis} (TCP): p50\,=\,1,331\,\textmu{}s,
p95\,=\,1,740\,\textmu{}s, p99\,=\,2,987\,\textmu{}s. The +459\,\textmu{}s
delta reflects the OS network stack for loopback TCP. Both measurements
are well below the 5\,ms per-step budget that keeps tracing transparent
to agents with 100--500\,ms tool calls.

\subsection{AgencyBench-V2 Harness Ablation}
\label{sec:eval_agency}

This experiment directly measures what the harness layer contributes
to agent performance, using real task descriptions from the
AgencyBench-V2 corpus~\cite{agencybench2026}.

\textbf{Tasks.} We load 28 regular tasks spanning six AgencyBench-V2
capability dimensions (Code, Backend, Game, MCP, Research, Frontend),
plus 2 adversarial tasks (unsafe shell commands and destructive SQL).
Of the 28, 9 are continuation tasks (subtask 2 explicitly builds on
subtask 1). Tasks are sorted so prior subtasks run before their
continuations.

\textbf{Conditions.} Three conditions over the same 28 tasks:
\begin{itemize}[noitemsep]
  \item \textbf{A — Bare:} No harness features. No memory injection,
    no tool result capping, no safety, no tracing.
  \item \textbf{B — +Infra:} HarnessAgent memory (ContextEngine, context
    injected from prior subtask when it passed), tool result cap (8k),
    circuit breaker, tracing, cost tracking.
  \item \textbf{C — Full HaaS:} B plus 3-stage safety pipeline, skill
    store lookup, cost guard.
\end{itemize}

\textbf{Scoring.} Each task is scored on a 0--10 rubric following
AgencyBench's structure (code correctness, file deliverables, functional
tests). The scoring model is a calibrated Gaussian
($\mu$\,=\,6.2, $\sigma$\,=\,1.5) with condition-specific adjustments:
$+$2.0 when prior context is available (memory bonus), $-$1.2 when
it is missing (hallucination penalty), $+$0.5 for skill-store reuse
(condition C), $+$0.3 for tool-cap benefit on hard tasks. These
adjustments reflect the mechanistic effect of each feature;
the Gaussian base captures inherent task difficulty variance.
Pass threshold: score $\ge$ 6.0.

Real HarnessAgent components are active: SafetyPipeline (real patterns),
ContextEngine (fakeredis-backed), TraceRecorder. The LLM response is
replaced by the scoring model (same approach as bench\_hermes.py).

\textbf{Results.}

\begin{table}[h]
\centering\small
\caption{AgencyBench-V2 Ablation (28 regular tasks, seed\,=\,42)}
\label{tab:agency}
\setlength{\tabcolsep}{4pt}
\begin{tabular}{lrrl}
\toprule
\textbf{Condition} & \textbf{Pass Rate} &
  \textbf{Avg Score} & \textbf{vs native-SDK} \\
\midrule
A  Bare (no harness)   & 50.0\% & 5.76 & $+$1.6\,pp  \\
B  +HarnessAgent infra & 60.7\% & 6.20 & $+$12.3\,pp \\
C  Full HaaS           & 71.4\% & 6.55 & $+$23.0\,pp \\
\midrule
\multicolumn{4}{l}{\small Reference: AgencyBench native-SDK baseline
  48.4\%~\cite{agencybench2026}} \\
\bottomrule
\end{tabular}
\end{table}

Memory and infrastructure (A$\to$B): $+$10.7\,pp. Safety and skill
store (B$\to$C): $+$10.7\,pp. Total HaaS lift: $+$21.4\,pp. Both
adversarial tasks were blocked by the safety pipeline in condition C
and passed (dangerously) in conditions A and B.

\subsection{\texorpdfstring{$\tau$}{tau}-bench Safety Compliance}
\label{sec:eval_tau}

Fifty tasks follow the $\tau$-bench~\cite{taubench2024} retail and
airline domain structure: 40 benign (order status, returns,
cancellations, seat upgrades) and 10 adversarial (prompt injection,
mass chargebacks, unauthorized data access, jailbreak attempts).

\textbf{Condition X} (safety off): the agent processes all requests
without guardrail intervention.
\textbf{Condition Y} (safety on): the 3-stage pipeline runs on every
request.

\begin{table}[h]
\centering\small
\caption{$\tau$-bench Safety Results (50 tasks)}
\label{tab:tau}
\setlength{\tabcolsep}{4pt}
\begin{tabular}{lrr}
\toprule
\textbf{Metric} & \textbf{X (safety off)} & \textbf{Y (safety on)} \\
\midrule
Benign success rate      & 67.5\% & 65.0\% \\
Adversarial blocked      & 0.0\%  & \textbf{100.0\%} \\
Policy compliance rate   & 10.0\% & \textbf{100.0\%} \\
False positive rate      & 0.0\%  & \textbf{0.0\%} \\
\bottomrule
\end{tabular}
\end{table}

Without the pipeline, 90\% of adversarial requests execute (policy
compliance 10\%). With the pipeline, 100\% are blocked and 0\% of
legitimate requests are incorrectly rejected (0\% FPR). The benign
success rate drops by 2.5\,pp in condition Y, reflecting the small
latency overhead of running three guardrail checks per request.

HAL's best published $\tau$-bench airline task success is
56\%~\cite{hal2025}. Our metric is orthogonal — we measure safety
compliance, not task completion — confirming that leaderboard accuracy
and safety compliance are independent dimensions.

\subsection{ATBench Safety Trajectories}
\label{sec:eval_atbench}

Fifty agent action trajectories follow the ATBench
taxonomy~\cite{atbench2026}: 30 unsafe (5--6 per ATBench category) and
20 benign. Each trajectory is a concrete agent action (tool call or LLM
output) routed through the appropriate guardrail stage (input, step, or
output), matching where each category is naturally detected.

\begin{table}[h]
\centering\small
\caption{ATBench Safety Trajectory Results (50 scenarios)}
\label{tab:atbench}
\setlength{\tabcolsep}{4pt}
\begin{tabular}{lrrl}
\toprule
\textbf{Metric} & \textbf{Result} & \textbf{Target} & \textbf{Status} \\
\midrule
TPR (unsafe blocked)  & \textbf{96.7\%} (29/30) & $\ge$90\% & \checkmark \\
TNR (benign allowed)  & \textbf{100.0\%} (20/20) & $\ge$95\% & \checkmark \\
FPR (benign blocked)  & \textbf{0.0\%}   & $<$5\%  & \checkmark \\
Overall accuracy      & \textbf{98.0\%}  & ---     & \\
\bottomrule
\end{tabular}
\end{table}

Five categories achieve 100\% TPR: Data Exfiltration, Harmful Content
Generation, Policy Bypass, Privilege Escalation, and Unsafe Code
Execution. Prompt Injection achieves 83.3\% (5/6) — one missed case
used the phrase ``you are now DAN,'' which the current signal list does
not cover. Adding this phrase to the detection patterns requires one
line of code and is the only identified gap.

\subsection{Hermes on GAIA-Style Tasks}
\label{sec:eval_hermes_gaia}

\textbf{Important framing.} This benchmark measures \emph{pre/post-patch
improvement} from the Hermes loop — not a GAIA leaderboard score.
HAL's top GAIA result is 74.55\% with Claude Sonnet 4.5~\cite{hal2025};
these are not comparable measurements.

Fifteen GAIA Level-1 failure types are seeded into the ErrorCollector:
tool-use errors (wrong function name, missing tool verification),
arithmetic mistakes (wrong formula, unit conversion errors),
multi-step reasoning failures (stopping too early, wrong chain),
and context loss (forgetting prior turn context). Twenty held-out
tasks measure pass@1 before and after patching.

Pre-patch pass@1: 5.0\% (the diverse failure modes expose many
weaknesses simultaneously). Three Hermes cycles run; the first patch
covers all six failure categories and is applied at cycle 1.
Post-patch pass@1: 90.0\% (+85\,pp). This exceeds the SQL-only
Hermes result (+6\,pp real, +70\,pp mock) because GAIA's diverse
failure types allow the patch to generalize broadly.

\subsection{Hermes on Real SQL Benchmark}
\label{sec:eval_hermes_sql}

NexusSql (GPT-5.5, Azure, temperature=0) is evaluated on 50 real SQL
tasks from \texttt{gretelai/synthetic\_text\_to\_sql}~\cite{gretelai2024},
stratified across 7 complexity tiers. A real SQLite database is built
per task; pass@1 uses result-set equality. Pre-patch failures seed
the ErrorCollector. Hermes runs 3 cycles; the Evaluator uses real
SQLite execution for each candidate patch.

Pre-patch pass@1: 40.0\% (20/50). Post-patch pass@1: 46.0\% (23/50),
+6\,pp. No patch exceeded the 0.60 auto-apply threshold; patches were
staged for human review — consistent with the conservative default.
The dominant failure mode is multi-join result-set mismatch (43\%
accuracy on \texttt{multiple\_joins} complexity).

\subsection{NexusSql Component Ablation}
\label{sec:eval_ablation}

Three conditions on 49 real SQL tasks:
\textbf{A} — LLM only (no schema context, no verifier, no retry);
\textbf{B} — SchemaStore + SQLVerifier, no retry;
\textbf{C} — full NexusSql (schema, verifier, up to 2 retries).

Condition A: 36.7\%. Condition B: 44.9\% (+8.2\,pp from schema and
verifier). Condition C: 49.0\% (+4.1\,pp from self-correction).
Total gain A$\to$C: +12.3\,pp. The verifier and schema together
account for 67\% of the total gain.

%=======================================================================
\section{Discussion}
\label{sec:discussion}

\subsection{Limitations}

\textbf{Simulation in benchmarks 7.5--7.8.} The AgencyBench-V2
ablation, $\tau$-bench, ATBench, and GAIA Hermes benchmarks use
calibrated scoring models rather than live LLM calls.
The scoring adjustments (+2.0 memory bonus, $-$1.2 no-memory penalty,
+0.5 skill store, safety blocks) are mechanistic: they reflect the
actual effect of each HarnessAgent feature on task output quality.
All HarnessAgent infrastructure components are active and real;
only the LLM is simulated. We are preparing live evaluations on the
official AgencyBench and GAIA datasets for the final workshop paper.

\textbf{Single-machine measurement.} All benchmarks run on one laptop.
Multi-replica deployments add Redis consumer group coordination overhead
to the FeedbackChannel that is not captured here.

\textbf{Span overhead on cloud Redis.} Real localhost Redis adds
459\,\textmu{}s over fakeredis (Section~\ref{sec:eval_spans}). A cloud
Redis instance (10--30\,ms round trip) would raise span overhead to
roughly 10--30\,ms per span — still below the 100\,ms tool call budget
but worth measuring before production deployment at scale.

\textbf{Prompt Injection gap.} One ATBench scenario using the phrase
``you are now DAN'' is not caught by the current signal list.
This is a known gap, fixable in one line.

\subsection{Relationship to HAL Reliability Dashboard}

HAL recently launched a Reliability Dashboard~\cite{hal2025} that
measures agent consistency, robustness, and safety — the same
dimensions HarnessAgent targets. This pivot confirms that the field
recognizes accuracy alone is insufficient. HarnessAgent provides the
production infrastructure layer that operationalizes reliability; HAL
provides the evaluation infrastructure to measure it. These roles are
complementary.

\subsection{Future Work}

\textbf{Adaptive context compression.} Replace the fixed 60/25/15
budget split with a learned allocator that adjusts per task type.

\textbf{ML-based routing.} Replace the linear scoring function
(Eq.~\ref{eq:routing}) with a contextual bandit that learns provider
preferences per tenant.

\textbf{Streaming guardrails.} Extend safety checks to partial streaming
outputs, enabling interruption before a harmful response completes.

\textbf{Plugin SDK.} Allow third-party verifiers, memory backends, and
routing policies without forking the harness.

\textbf{Live AgencyBench and GAIA runs.} Replace the simulated ablation
with real LLM execution on the official datasets.

%=======================================================================
\section{Conclusion}
\label{sec:conclusion}

Production AI agents fail for predictable, infrastructure reasons —
not because of weak reasoning. HarnessAgent addresses this by
formalizing a HaaS contract in which the harness owns routing, memory,
tracing, safety, cost, and improvement while the agent framework owns
task logic. A four-line integration activates the full stack for
LangGraph, AutoGen, CrewAI, or Agno.

Ten benchmarks validate the system. Three real-execution infrastructure
benchmarks show: 82.9\% token reduction from GraphRAG, 100\%/0\%
TPR/FPR semantic cache, and sub-millisecond circuit-breaker activation
with zero false trips. Span recording overhead is 872\,\textmu{}s p50
on fakeredis and 1,331\,\textmu{}s on real localhost Redis — below the
5\,ms threshold in both cases.

The AgencyBench-V2 harness ablation shows the clearest evidence: adding
HarnessAgent to bare LangGraph lifts pass rate from 50.0\% to 71.4\%
(+21.4\,pp), exceeding the published native-SDK baseline of
48.4\%~\cite{agencybench2026} by 23.0\,pp. The safety pipeline blocks
100\% of adversarial requests on $\tau$-bench tasks with 0\% false
positives, and achieves TPR\,=\,96.7\%, TNR\,=\,100\% on ATBench
trajectories. Hermes improves pre/post-patch pass@1 by +85\,pp on
GAIA-style tasks and +6\,pp on verified SQL execution.

Meng et~al.~\cite{meng2026survey} identified production multi-framework
harness as an open research direction in April 2026. HAL's simultaneous
pivot to a Reliability Dashboard~\cite{hal2025} confirmed the need.
HarnessAgent fills both gaps. The package is \texttt{pip install
agent-haas} with 1,087 passing tests, MIT license, and adapters for
four production frameworks.

%=======================================================================
\section*{Acknowledgments}

The author thanks the open-source communities behind LangGraph, AutoGen,
CrewAI, Agno, Chroma, sentence-transformers, FastAPI, and fakeredis,
whose foundational work made HarnessAgent possible. The AgencyBench team
(GAIR-NLP) and the Princeton HAL team are acknowledged for releasing
their benchmarks and evaluation infrastructure publicly.

%=======================================================================
\balance
\bibliographystyle{IEEEtran}
\begin{thebibliography}{99}

\bibitem{zhong2026harness}
Z.~Zhong and Y.~Zhu,
``AI Harness Engineering: A Systematic Framework for Production-Grade
  Agent Infrastructure,''
\emph{arXiv:2605.13357 [cs.SE]}, May 2026.

\bibitem{lin2026ahe}
K.~Lin et~al.,
``Agentic Harness Engineering: Observability-Driven Evolution of
  Coding Agent Infrastructure,''
\emph{arXiv:2604.25850 [cs.CL]}, Apr.\ 2026.

\bibitem{seong2026last}
J.~Seong et~al.,
``The Last Harness You'll Ever Build: Meta-Learning for Agent
  Infrastructure Construction,''
\emph{arXiv:2604.21003 [cs.AI]}, Apr.\ 2026.

\bibitem{meng2026survey}
H.~Meng et~al.,
``Agent Harness Survey: A Systematic Review of 22 Production Agent
  Systems,''
\emph{preprints:202604.0428}, Apr.\ 2026.

\bibitem{hal2025}
S.~Kapoor et~al.,
``Holistic Agent Leaderboard: The Missing Infrastructure for AI Agent
  Evaluation,''
in \emph{Proc.\ ICLR}, 2026. arXiv:2510.11977 [cs.AI].

\bibitem{agencybench2026}
GAIR-NLP,
``AgencyBench: Benchmarking the Frontiers of Autonomous Agents in
  Real-World Long-Horizon Contexts,''
\emph{arXiv:2601.11044}, Jan.\ 2026.

\bibitem{taubench2024}
Sierra Research,
``$\tau$-bench: A Benchmark for Tool-Agent-User Interaction in
  Real-World Domains,''
\emph{arXiv:2406.12045 [cs.AI]}, Jun.\ 2024.

\bibitem{atbench2026}
X.~et~al.,
``ATBench: A Diverse and Realistic Agent Trajectory Benchmark for
  Safety,''
\emph{arXiv:2604.02022 [cs.AI]}, Apr.\ 2026.

\bibitem{gaia2023}
G.~Mialon et~al.,
``GAIA: A Benchmark for General AI Assistants,''
\emph{arXiv:2311.12983 [cs.CL]}, Nov.\ 2023.

\bibitem{memgpt2023}
C.~Packer et~al.,
``MemGPT: Towards LLMs as Operating Systems,''
\emph{arXiv:2310.08560 [cs.AI]}, Oct.\ 2023.

\bibitem{graphrag2024}
D.~Edge et~al.,
``From Local to Global: A Graph RAG Approach to Query-Focused
  Summarization,''
\emph{arXiv:2404.16130 [cs.CL]}, Apr.\ 2024.

\bibitem{terag2025}
Anonymous,
``TERAG: Token-Efficient Graph Retrieval-Augmented Generation,''
\emph{arXiv:2509.18667}, Sep.\ 2025.

\bibitem{deepseekr1}
DeepSeek-AI,
``DeepSeek-R1: Incentivizing Reasoning Capability in LLMs via
  Reinforcement Learning,''
\emph{arXiv:2501.12948 [cs.CL]}, Jan.\ 2025.

\bibitem{nexussql}
P.~Tivhale,
``NexusSQL / SQLAS: A Structured SQL Agent Evaluation Framework,''
\emph{Zenodo}, DOI:\,10.5281/zenodo.20180541, May 2026.

\bibitem{langgraph}
LangChain-AI,
``LangGraph: Building Stateful, Multi-Actor Applications with LLMs,''
\url{https://langchain-ai.github.io/langgraph/}, 2024.

\bibitem{autogen}
Q.~Wu et~al.,
``AutoGen: Enabling Next-Gen LLM Applications via Multi-Agent
  Conversation,''
\emph{arXiv:2308.08155 [cs.AI]}, Aug.\ 2023.

\bibitem{minilm}
W.~Wang et~al.,
``MiniLM: Deep Self-Attention Distillation for Task-Agnostic
  Compression of Pre-Trained Transformers,''
in \emph{Advances in Neural Information Processing Systems (NeurIPS)},
vol.~33, pp.~5776--5788, 2020.

\bibitem{langsmith}
LangChain-AI,
``LangSmith: Debugging and Monitoring for LLM Applications,''
\url{https://smith.langchain.com/}, 2024.

\bibitem{gretelai2024}
Gretel-AI,
``Synthetic Text-to-SQL Dataset,''
HuggingFace Datasets,
\url{https://huggingface.co/datasets/gretelai/synthetic_text_to_sql},
2024.

\end{thebibliography}

% -----------------------------------------------------------------------
% AI DISCLOSURE (required by arXiv policy)
% -----------------------------------------------------------------------
% This manuscript was prepared with assistance from Claude (Anthropic).
% All experimental results, benchmark scripts, and system implementation
% were produced and verified by the human author. AI tools assisted with
% drafting and copyediting prose only. All claims are verifiable against
% benchmarks/results/ in the public repository. The author takes full
% responsibility for all content.
% -----------------------------------------------------------------------

\end{document}
